% Ad Familiares 13.60 --- headnote
% ad-familiares-13-60, probably late 46 BC, Rome

\textit{Cicero at Rome to Gaius Munatius, son of Gaius, written
after his return from the eastern campaign, probably late in 46
BC (the manuscript dateline: \textit{Scr. Romae post reditum
fortasse ex. a. 708 (46)}, where the Perseus text's
\textit{ex. a. 698} appears to be a transcription slip for
\textit{ex. a. 708}). A short recommendation of Lucius
Livineius Trypho, freedman of Cicero's friend L. Regulus, whose
own \textit{calamitas} --- presumably his civil-war
condemnation --- has put Cicero under a renewed obligation to
his household. The freedman, however, is loved on his own
merits, for services rendered during ``those times of ours''
(the civil-war years following Pharsalus, when Cicero was
detained at Brundisium): Trypho had run many risks for Cicero's
sake and made repeated voyages through the depth of winter --- a
glancing reference to the courier duties freedmen of well-placed
households so often performed for their patrons' friends in
that period.}
